\section*{Chapitre \ref{ch:b3:renal-osmoregulation} --- Physiologie rénale et osmorégulation}

\begin{solution}{exo:b3:renal-osmoregulation:1}
\omterm{def:b3:renal-osmoregulation:nephron}{Glomérule}~: filtre le plasma moins les protéines. \omterm{def:b3:renal-osmoregulation:nephron}{Tube proximal}~:
réabsorbe les deux tiers du sel et de l'eau, tout le glucose et les
acides aminés, le bicarbonate. Branche descendante de Henle~: perd son
eau vers la médullaire. Branche ascendante~: pompe le NaCl dehors,
imperméable à l'eau. Tube distal~: ajuste le sodium (\omterm{def:b3:renal-osmoregulation:controls}{aldostérone}),
sécrète le potassium. \omterm{def:b3:renal-osmoregulation:nephron}{Tube collecteur}~: réabsorption d'eau sous
vasopressine, traitement des protons et de l'urée.
\end{solution}

\begin{solution}{exo:b3:renal-osmoregulation:2}
La clairance~: le volume de plasma épuré de la substance par unité de
temps, $UV/P$. Ici $100\times 1.5/2 = \qty{75}{mL/min}$, au-dessous du
DFG de \qty{120}{mL/min}~: la substance est filtrée et partiellement
réabsorbée (environ \qty{38}{\%} de la quantité filtrée).
\end{solution}

\begin{solution}{exo:b3:renal-osmoregulation:3}
Un poisson d'eau douce est plus salé que son milieu~: l'eau entre par
les branchies par osmose, si bien que boire ne ferait qu'ajouter au
flot~; il excrète une urine abondante et diluée et les \omterm{def:b3:renal-osmoregulation:comparative}{cellules à
chlorure} de ses branchies captent activement Na$^{+}$ et Cl$^{-}$ dans
l'eau. Un poisson marin est moins salé que la mer~: il perd de l'eau par
les branchies, boit de l'eau de mer pour la remplacer, et les \omterm{def:b3:renal-osmoregulation:comparative}{cellules à
chlorure} rejettent le sel avalé.
\end{solution}

\begin{solution}{exo:b3:renal-osmoregulation:4}
L'\omterm{prop:b3:renal-osmoregulation:countercurrent}{effet unitaire} est la différence de \qty{200}{mOsm/L} que la branche
ascendante large maintient entre l'interstitium et son propre liquide en
pompant le NaCl dehors sans laisser l'eau suivre. La branche ascendante
doit être imperméable à l'eau, sans quoi l'eau suivrait le sel et la
différence s'évanouirait~; la branche descendante doit être perméable à
l'eau pour que le liquide qui atteint le coude ait déjà été concentré
par l'interstitium et que chaque étage travaille sur un liquide plus
salé que le précédent. Deux branches de même perméabilité ne pourraient
pas multiplier.
\end{solution}

\begin{solution}{exo:b3:renal-osmoregulation:5}
Extrémité afférente~: $60 - 18 - 28 = \qty{14}{mmHg}$~; extrémité
efférente~: $60 - 18 -
35 = \qty{7}{mmHg}$. À mesure que l'eau est filtrée les protéines
laissées derrière se concentrent et la pression oncotique monte le long
du capillaire~; si elle atteint \qty{42}{mmHg} la pression nette est
nulle et la filtration s'arrête avant la fin (l'équilibre de
filtration). Le DFG dépend alors du débit plasmatique~: un débit plus
rapide concentre les protéines plus lentement, si bien que la filtration
se poursuit sur une plus grande longueur et que le DFG monte avec le
débit.
\end{solution}

\begin{solution}{exo:b3:renal-osmoregulation:6}
Charge filtrée~: $\qty{40}{mL/min}\times\qty{0.14}{mmol/mL} =
\qty{5.6}{mmol/min}$. Excrétion~: $70\times 1 = \qty{0.07}{mmol/min}$.
Réabsorbé~: $5.53/5.6 = \qty{98.75}{\%}$. Clairance du sodium~: $70\times
1/140 = \qty{0.5}{mL/min}$.
\end{solution}

\begin{solution}{exo:b3:renal-osmoregulation:7}
$900/1200 = \qty{0.75}{L}$ par jour. À \qty{400}{mOsm/L}~: $900/400 =
\qty{2.25}{L}$ d'urine, donc environ $2.25 + 0.9 = \qty{3.15}{L}$ à
boire (moins l'eau des aliments)~: le patient dont le mécanisme de
concentration défaille doit boire, et il est en danger s'il ne le peut
pas.
\end{solution}

\begin{solution}{exo:b3:renal-osmoregulation:8}
(a) Diabète insipide~: \qtyrange{10}{15}{L} d'urine à
\qtyrange{50}{100}{mOsm/L}, sodium plasmatique élevé dès que la boisson
tarde, soif implacable. (b) Vasopressine inappropriée~: urine rare et
concentrée, eau retenue, sodium plasmatique bas par dilution, soif non
accrue bien que la boisson continue --- confusion et convulsions quand
le sodium tombe bas. (c) Excès d'\omterm{def:b3:renal-osmoregulation:controls}{aldostérone}~: sodium retenu avec
l'eau, volume augmenté, pression artérielle élevée, potassium perdu
(K$^{+}$ plasmatique bas), sodium plasmatique presque normal puisque
l'eau suit le sel et que le rein échappe en partie~; volume urinaire
normal. (d) Régime sans sel~: rénine et \omterm{def:b3:renal-osmoregulation:controls}{aldostérone} montent en un jour,
le sodium urinaire tombe à quelques millimoles par jour, le sodium
plasmatique est tenu, le volume et la pression sont un peu plus bas.
\end{solution}

\begin{solution}{exo:b3:renal-osmoregulation:9}
Urine diluée~: clairance osmolaire $100\times 6/300 = \qty{2}{mL/min}$~;
\omterm{prop:b3:renal-osmoregulation:water}{clairance de l'eau libre} $6 - 2 = +\qty{4}{mL/min}$~: le rein se défait
d'eau sans soluté, comme après une charge d'eau, vasopressine éteinte.
Urine concentrée~: clairance osmolaire $1000\times 0.6/300 =
\qty{2}{mL/min}$~; \omterm{prop:b3:renal-osmoregulation:water}{clairance de l'eau libre} $0.6 - 2 = -\qty{1.4}{mL/min}$~:
\qty{1.4}{mL/min} d'eau libre sont rendus au corps --- déshydratation,
vasopressine élevée.
\end{solution}

\begin{solution}{exo:b3:renal-osmoregulation:10}
Numérotons les étages depuis le cortex~; à l'étage $k$, notons
l'interstitium $I_{k}$. Le liquide descendant en $k$ vaut $I_{k}$~; le liquide
ascendant en $k$ vaut $I_{k} - 200$~; et le liquide ascendant en $k$ est
celui qui a quitté le coude et franchi les étages $n, n-1, \dots, k+1$.
En régime stationnaire le liquide entrant à \qty{300}{} quitte le haut
de la branche ascendante à $I_{1} - 200$, et l'interstitium de chaque
étage dépasse celui du dessus d'au plus l'\omterm{prop:b3:renal-osmoregulation:countercurrent}{effet unitaire}, donc
$I_{k} \le 300 + 200k$ avec égalité quand chaque étage travaille à plein
effet~: le coude atteint $300 + 200n$. Le gradient est donc borné par la
longueur de l'anse (le nombre d'étages) multipliée par l'\omterm{prop:b3:renal-osmoregulation:countercurrent}{effet unitaire},
et non par la puissance de pompage d'une cellule. $n$ est la longueur de
l'anse divisée par la distance sur laquelle le liquide s'équilibre avec
l'interstitium~: les longues anses juxtamédullaires, et les très longues
anses des rongeurs du désert, donnent un grand $n$~; la capacité de la
pompe et le lessivage du soluté par les vasa recta plafonnent l'effet.
\end{solution}

\begin{solution}{exo:b3:renal-osmoregulation:11}
\omterm{def:b3:renal-osmoregulation:comparative}{Eau métabolique}~: $5\times 0.8 = \qty{4}{g}$ d'amidon donnent $4\times 0.6
= \qty{2.4}{g}$. Pertes~: urine $3/5500\ \unit{L} = \qty{0.55}{mL}$~;
évaporation \qty{1.5}{g} (déjà net de la récupération nasale). Perte
totale \qty{2.05}{g} contre un gain de \qty{2.4}{g}~: un excédent de
\qty{0.35}{g}, qui couvre la faible perte fécale, et l'eau hygroscopique
des graines est un bonus. Il survit sans boire. Sans la récupération
nasale la perte respiratoire serait de $1.5/0.3 = \qty{5}{g}$, et il
mourrait en quelques jours.
\end{solution}

\begin{solution}{exo:b3:renal-osmoregulation:12}
Dans un écoulement parallèle le sang et le dialysat s'équilibrent à
mi-chemin le long de la membrane et la différence de concentration qui
fait diffuser tombe à zéro~: au mieux le sang repart avec la moitié de
son urée. À contre-courant le dialysat le plus frais rencontre le sang
le plus épuré et le gradient est maintenu sur toute la membrane, de
sorte que le sang peut repartir presque à la concentration d'entrée du
dialysat --- le principe même de l'\omterm{def:b3:renal-osmoregulation:nephron}{anse de Henle}. Clairance~: $300\times 0.7 =
\qty{210}{mL/min}$ pendant $3\times 4 = \qty{12}{h}$ par semaine, soit
$210\times 720 = \qty{151}{L}$ par semaine~; deux reins à
\qty{125}{mL/min} épurent $125\times\num{10080} = \qty{1260}{L}$. La
machine fait environ le huitième du travail des reins, soit
\qty{15}{mL/min} en moyenne --- de quoi vivre, pas de quoi se bien
porter.
\end{solution}

\begin{solution}{pb:b3:renal-osmoregulation:1}
\textbf{1.} Pression nette $55 - 15 - 30 = \qty{10}{mmHg}$~; DFG $= 12.5
\times 10 = \qty{125}{mL/min}$.
\textbf{2.} $125\times 1440 = \qty{180}{L/d}$~; $180/3 = 60$ fois.
\textbf{3.} Inuline~: $62.5\times 1/0.5 = \qty{125}{mL/min}$, égale au
DFG. Créatinine~: $1.3\times 1/0.01 = \qty{130}{mL/min}$, un peu plus
grande parce que le \omterm{def:b3:renal-osmoregulation:nephron}{tube proximal} sécrète un peu de créatinine en plus
de ce qui est filtré.
\textbf{4.} Clairance de l'APAH $13\times 1/0.02 = \qty{650}{mL/min}$, le
débit plasmatique rénal~; débit sanguin $650/(1 - 0.45) = \qty{1180}{mL/min}$,
soit environ \qty{1.2}{L/min}~; fraction de filtration $125/650 = 0.19$.
\textbf{5.} Pression nette \qty{15}{mmHg}, DFG \qty{188}{mL/min} (moins
en réalité, puisque la pression oncotique monte le long du capillaire).
L'angiotensine~II resserre l'artériole efférente et élève la pression
glomérulaire~; la bloquer abaisse un peu la pression et le DFG. Dans le
diabète les \omterm{def:b3:renal-osmoregulation:nephron}{glomérules} filtrent sous haute pression et le filtre s'use~;
abaisser la pression retarde l'atteinte de plusieurs années.
\textbf{6.} Les podocytes et leurs diaphragmes de fente (ou la charge de
la membrane basale)~: la couche qui retient l'albumine. Perdre
\qty{3}{g/d} abaisse la pression oncotique plasmatique~; le liquide
quitte les capillaires partout et les tissus enflent --- l'œdème du
syndrome néphrotique.
\textbf{7.} Filtré $125\times 5 = \qty{0.625}{mmol/min}$~; excrété
$250\times 1 = \qty{0.25}{mmol/min}$~; réabsorbé $0.375$, soit
\qty{60}{\%}~; clairance $250/5 = \qty{50}{mL/min}$.
\textbf{8.} Filtré $180\times 0.14 = \qty{25.2}{mol/d}$ de Na$^{+}$,
\qty{580}{g} de sodium, \qty{1.47}{kg} en NaCl. Excrété $100\times 1.44
= \qty{144}{mmol/d}$~; réabsorbé $1 - 144/\num{25200} = \qty{99.4}{\%}$.
Près d'un kilogramme et demi de sel serait perdu par jour.
\textbf{9.} Filtré $125\times 5 = \qty{0.625}{mmol/min}$, soit
\qty{112}{mg/min}, \qty{162}{g/d}~; rien n'est excrété, puisqu'on est
au-dessous de $T_{m}$. La charge filtrée atteint $T_{m}$ à $2.1/0.125 =
\qty{16.8}{mmol/L}$ (\qty{300}{mg/dL}) --- plus haut que le seuil observé
d'environ \qty{10}{mmol/L}, parce que les \omterm{def:b3:renal-osmoregulation:nephron}{néphrons} diffèrent et que les
plus faibles débordent les premiers (l'étalement de la courbe de
titration).
\textbf{10.} Filtré $125\times 20 = \qty{2.5}{mmol/min}$~; excrété
$2.5 - 2.1 = \qty{0.4}{mmol/min}$, \qty{72}{mg/min}, soit $576$ mmol ou
\qty{104}{g} par jour. Urine supplémentaire $576/300 = \qty{1.9}{L}$. Le
glucose entraîne l'eau dans l'urine (diurèse osmotique), le patient se
déshydrate et a soif~; \qty{104}{g} de glucose font \qty{1.8}{MJ} perdus
par jour, et les cellules, incapables de capter le glucose, brûlent
graisses et protéines~: le poids tombe.
\textbf{11.} Réabsorbé $\approx \qty{25}{mol/d}$~; ATP $25/3 =
\qty{8.4}{mol/d}$~; $8.4\times 50 = \qty{420}{kJ}$, environ
\qty{6}{\%} de \qty{7}{MJ}.
\textbf{12.} Pour sécréter chaque déchet le tube aurait besoin d'un
transporteur qui le reconnaisse --- des milliers d'entre eux, et aucun
pour une molécule jamais rencontrée. Tout filtrer de ce qui est petit et
réabsorber les quelques centaines d'espèces que le corps veut ne demande
de transporteurs que pour ce qu'on sait précieux~; tout ce qui n'est pas
reconnu s'en va par défaut. Le prix est l'énergie de la question~11.
\textbf{13.} Clairance osmolaire $600\times 1/290 = \qty{2.07}{mL/min}$~;
\omterm{prop:b3:renal-osmoregulation:water}{clairance de l'eau libre} $1 - 2.07 = -\qty{1.07}{mL/min}$~: le rein
conserve l'eau.
\textbf{14.} Minimum $600/1200 = \qty{0.5}{L}$~; maximum $600/50 =
\qty{12}{L}$~; besoin quotidien au minimum $0.5 + 0.9 = \qty{1.4}{L}$.
\textbf{15.} \qty{1000}{mOsm} à \qty{1200}{mOsm/L} demandent
\qty{0.83}{L} d'urine~: gain net \qty{0.17}{L} avant le soluté du jour.
Après~: le soluté total du jour fait \qty{1600}{mOsm}, soit \qty{1.33}{L}
d'urine, plus \qty{0.9}{L} d'insensible, \qty{2.23}{L} dehors pour
\qty{1}{L} dedans --- un déficit de \qty{1.23}{L}, contre \qty{1.4}{L}
sans rien boire. Le litre d'eau de mer n'améliore le bilan que de
\qty{0.17}{L}.
\textbf{16.} Urine à \qty{50}{mOsm/L} pour la journée~: $600/50 =
\qty{12}{L}$~; il faut boire \qty{12.9}{L} par jour, soit environ \qty{13}{L}.
\textbf{17.} $B$ litres de bière apportent $20B$ mOsm~; l'urine maximale
vaut $(600 + 20B)/50 = 12 + 0.4B$ litres, et elle doit emporter $B - 0.9$
litres~: $B - 0.9 \le 12 + 0.4B$, $B \le \qty{21.5}{L}$. Une vingtaine
de litres --- mais si l'on ne mange presque rien le soluté du jour tombe
à \qty{200}{mOsm} et la limite descend vers \qty{8}{L}~:
l'hyponatrémie des gros buveurs qui ne mangent pas.
\textbf{18.} $298/290 - 1 = \qty{2.8}{\%}$~: bien au-delà du \qty{1}{\%}
qui déclenche la vasopressine et la soif~; toutes deux sont maximales.
À soluté constant~: $290\times 42 = 298\times V$, $V = \qty{40.9}{L}$, un
déficit d'environ \qty{1.1}{L}.
\textbf{19.} Six litres d'eau diluent le plasma tandis que la sueur en
retire le sel~; et l'exercice, la douleur et le stress libèrent de la
vasopressine quelle que soit l'\omterm{def:b3:renal-osmoregulation:problem}{osmolarité}, de sorte que le rein ne peut
diluer l'urine assez pour excréter l'excès. Le sodium plasmatique tombe,
les cellules du cerveau gonflent --- convulsions et décès s'en sont
suivis. Il faudrait au rein une vasopressine éteinte et \qty{6}{L}
d'urine à $U_{\min}$~; le coureur devrait boire selon sa soif et
remplacer le sel.
\textbf{20.} $3/5500\ \unit{L} = \qty{0.55}{mL}$ par jour. Un rein humain
aurait besoin de $3/1200 = \qty{2.5}{mL}$ pour le même soluté~: le rat
kangourou emploie un cinquième de l'eau.
\textbf{21.} Entrées~: \qty{2.4}{g}. Sorties~: $1.6 + 0.3 + 0.55 = \qty{2.45}{g}$.
Équilibré à \qty{0.05}{g} près par jour, ce que fournit l'eau
hygroscopique des graines (quelques pour cent de \qty{5}{g}).
\textbf{22.} Eau $500\times 0.65 = \qty{325}{L}$~; le quart fait
\qty{81}{L}~; à \qty{5}{L/d}, seize jours.
\textbf{23.} $150/800 = \qty{0.19}{L}$, soit environ \qty{190}{mL} de
saumure --- moins que les \qty{100}{mL} d'eau de mer et l'eau du poisson
réunies, de sorte que l'oiseau gagne de l'eau. Son rein, à
\qty{600}{mOsm/L} (\qty{300}{mmol/L} de NaCl), aurait besoin de
\qty{0.5}{L} d'urine pour emporter le même sel, plus d'eau qu'il n'en a
pris~: un rein qui ne sait pas battre l'eau de mer ne peut laisser un
oiseau la boire, et la glande, au double de la concentration de la mer,
le peut.
\textbf{24.} \qty{30}{g} d'eau entrent, donc environ \qty{30}{mL} d'urine
par jour, à une petite fraction de l'\omterm{def:b3:renal-osmoregulation:problem}{osmolarité} du sang (quelques
dizaines de mOsm/L contre \num{300}). Le sel part dans cette urine et par
diffusion à travers les branchies~; sans capture active par les \omterm{def:b3:renal-osmoregulation:comparative}{cellules
à chlorure} dans une eau qui tient une millimole par litre, le poisson
serait épuisé en quelques jours.
\textbf{25.} DFG \qty{125}{mL/min}, \qty{180}{L} par jour~; urine
obligatoire \qty{0.5}{L} par jour~; un litre d'eau de mer rapporte un net
de \qty{0.17}{L} une fois son sel excrété --- presque rien.
\end{solution}
